% =====================================================================
\chapter{Datos, ventanas y objetivo de entrenamiento}
\label{tema:1}
% =====================================================================

Un modelo de lenguaje autorregresivo se entrena con una única secuencia larga de enteros, de la que se extraen ventanas de longitud fija. Este tema formaliza esa preparación de los datos, define la función que se minimiza (la entropía cruzada media sobre todas las posiciones de un lote) y demuestra tres propiedades de esa función que se usan durante todo el módulo: su gradiente respecto de los logits, su invariancia frente a traslaciones de los logits, que permite calcularla de forma estable, y su valor cuando la predicción es uniforme. Se sigue el §4.1 y el §5.1 de la tarea~1 de CS336~\cite{cs336}, el §9.3.3 de~\cite{d2l} y la clase~7 de CS224N~\cite{cs224n}.

\section{El corpus como secuencia de tokens}

\begin{definicion}[Corpus tokenizado]\label{def:corpus}
Sea $\V$ un vocabulario de tamaño $K$, identificado con $\{0,1,\dots,K-1\}$. Un \emph{corpus tokenizado} es una secuencia $u_{1:n}=(u_1,\dots,u_n)\in\V^{n}$. Si el texto original está formado por varios documentos, se concatenan en una sola secuencia, separados por un token delimitador.
\end{definicion}

La concatenación es una decisión de ingeniería con una justificación precisa: todas las ventanas tienen la misma longitud y no hace falta rellenar con símbolos de relleno (\emph{padding}), que consumen cómputo sin aportar señal~\cite[§5.1]{cs336}. El precio es que una ventana puede cruzar la frontera entre dos documentos; el delimitador permite al modelo aprender que lo que sigue no depende de lo anterior.

\addtocounter{definicion}{1}%

Un corpus sintético tiene una ventaja metodológica: se conoce el proceso que lo genera. Se sabe, por ejemplo, que tras «receta.» viene uno de seis verbos y que «un reloj» exige «negro» y no «negra». Esto permite juzgar qué ha aprendido el modelo y calcular la pérdida mínima alcanzable (\cref{tema:4}).

\section{Ventanas de contexto y lotes}

\begin{definicion}[Ventana de contexto]\label{def:ventana}
Sea $T\ge 1$ la \emph{longitud de contexto}. Para cada $i\in\{1,\dots,n-T\}$, la \emph{ventana} que empieza en $i$ es el par
\[
x=(u_i,u_{i+1},\dots,u_{i+T-1}),\qquad y=(u_{i+1},u_{i+2},\dots,u_{i+T}),
\]
es decir, $y_t=x_{t+1}$ para $t<T$ e $y_T=u_{i+T}$. La secuencia $y$ es la \emph{ventana objetivo}.
\end{definicion}

\begin{proposicion}\label{prop:ventanas}
Un corpus de $n$ tokens tiene exactamente $n-T$ ventanas de longitud $T$, y cada una contiene $T$ problemas de predicción: predecir $y_t$ a partir del prefijo $x_{1:t}$, para $t=1,\dots,T$.
\end{proposicion}

\addtocounter{definicion}{1}%

\begin{definicion}[Lote]\label{def:lote}
Un \emph{lote} de tamaño $B$ es un par de matrices de enteros $X,Y\in\V^{B\times T}$ cuyas filas son $B$ ventanas, con inicios elegidos de forma independiente y uniforme en $\{1,\dots,n-T\}$. Un lote contiene $B\cdot T$ predicciones.
\end{definicion}

Con $B=T=32$, cada paso de entrenamiento ve $1024$ predicciones; en $800$ pasos son $819\,200$, unas $57$ veces los $14\,459$ tokens de entrenamiento del \cref{ej:corpus}. Con tan pocos datos es inevitable repetir; los grandes modelos, en cambio, ven cada token una vez o pocas veces (módulo~16).

\begin{observacion}[Partición en entrenamiento y validación]\label{obs:particion}
El conjunto de validación se obtiene cortando el corpus \emph{por posición}: el 90~\% inicial ($14\,459$ tokens) se usa para entrenar y el 10~\% final ($1\,607$ tokens) para validar. Si se eligieran ventanas al azar para cada conjunto, una ventana de validación podría solaparse con una de entrenamiento y la pérdida de validación subestimaría el error de generalización.
\end{observacion}

La alternativa al muestreo independiente de inicios es partir el corpus en trozos consecutivos sin solape, con un desplazamiento inicial aleatorio en cada época, como hace~\cite[§9.3.3]{d2l}. El \cref{eje:1-d2l} compara ambas estrategias.

\section{La pérdida de entropía cruzada}

En cada posición $t$ de cada ventana, el modelo devuelve un vector de \emph{logits} $z_t\in\R^{K}$ y la distribución predicha $p_t=\softmax(z_t)$, con $p_{t,j}=e^{z_{t,j}}/\sum_{a}e^{z_{t,a}}$.

\begin{definicion}[Pérdida de entrenamiento]\label{def:perdida}
La pérdida de una posición con objetivo $y\in\V$ y logits $z$ es
\[
\ell(z,y)=-\ln\softmax(z)_y.
\]
La pérdida de un lote es la media sobre sus $B\cdot T$ posiciones,
\[
\mathcal{L}(\theta)=\frac{1}{BT}\sum_{b=1}^{B}\sum_{t=1}^{T}\ell\bigl(z_{b,t}(\theta),\,Y_{b,t}\bigr),
\]
donde $z_{b,t}(\theta)$ son los logits que el modelo de parámetros $\theta$ produce en la posición $t$ de la fila $b$ a partir de $X_{b,1:t}$.
\end{definicion}

La pérdida $\ell$ es la entropía cruzada entre la distribución degenerada en $y$ y la distribución predicha~\cite[§4.1]{cs336}. Su media sobre el corpus es la estimación de máxima verosimilitud del módulo~1 aplicada a un modelo con parámetros compartidos: minimizar $\mathcal{L}$ equivale a maximizar $\sum\ln P_\theta(y_t\mid x_{1:t})$.

\begin{definicion}[Perplejidad]\label{def:ppl}
Si una secuencia de $m$ predicciones tiene pérdidas $\ell_1,\dots,\ell_m$, su \emph{perplejidad} es
\[
\PPL=\exp\Bigl(\frac1m\sum_{i=1}^{m}\ell_i\Bigr).
\]
\end{definicion}

La perplejidad es el número de opciones equiprobables entre las que el modelo duda en media: un modelo uniforme sobre $K$ símbolos tiene perplejidad $K$. Es independiente de la base del logaritmo, porque $e^{\mathcal{L}}=2^{\mathcal{L}/\ln 2}$. En la práctica del módulo, el modelo de bigramas del módulo~1 llega a $2{,}035$ nats por carácter ($\PPL=7{,}65$) y el mini-GPT a $0{,}324$ ($\PPL=1{,}38$).

\section{Gradiente y estabilidad numérica}

\begin{proposicion}[Gradiente de la entropía cruzada]\label{prop:gradiente}
Sea $p=\softmax(z)$. Para todo $i\in\V$,
\[
\frac{\partial\ell(z,y)}{\partial z_i}=p_i-\ind[i=y],
\]
es decir, $\nabla_z\ell=p-e_y$, donde $e_y$ es el vector \emph{one-hot} de $y$.
\end{proposicion}

\addtocounter{definicion}{1}%

El gradiente de la \cref{prop:gradiente} es la única señal que recibe todo el mini-GPT: la retropropagación (módulo~10) la transporta desde los logits hasta los embeddings.

\begin{proposicion}[Invariancia frente a traslaciones]\label{prop:traslacion}
Para todo $c\in\R$, $\softmax(z+c\mathbf 1)=\softmax(z)$. En consecuencia, con $z^{\star}=\max_a z_a$,
\[
\ell(z,y)=-(z_y-z^{\star})+\ln\sum_{a}e^{\,z_a-z^{\star}},
\]
y todos los exponentes del segundo miembro son menores o iguales que~0.
\end{proposicion}

La fórmula estable evita el desbordamiento: en precisión simple el mayor número representable es unos $3{,}4\cdot 10^{38}\approx e^{88{,}7}$, de modo que $e^{1000}$ daría infinito y el cociente, un valor indefinido. Es la técnica \emph{log-sum-exp} del módulo~11, que CS336 exige en su implementación de la entropía cruzada~\cite[§4.1]{cs336}.

\begin{proposicion}[Pérdida de la predicción uniforme]\label{prop:lnK}
Si todos los logits de una posición son iguales, la pérdida vale $\ln K$ para cualquier objetivo.
\end{proposicion}

Con $K=46$, $\ln 46=3{,}829$ nats ($5{,}52$ bits). Una red recién inicializada con pesos pequeños produce logits casi iguales, y su pérdida inicial está cerca de ese valor; en la práctica del módulo el primer lote da $4{,}046$. Una pérdida inicial muy superior a $\ln K$ indica una inicialización defectuosa: el modelo está muy seguro de respuestas equivocadas.

